% This is samplepaper.tex, a sample chapter demonstrating the
% LLNCS macro package for Springer Computer Science proceedings;
% Version 2.21 of 2022/01/12
%
\documentclass{llncs}
%
\usepackage[T1]{fontenc}
% T1 fonts will be used to generate the final print and online PDFs,
% so please use T1 fonts in your manuscript whenever possible.
% Other font encondings may result in incorrect characters.
%
\usepackage{graphicx}
% Used for displaying a sample figure. If possible, figure files should
% be included in EPS format.
%
% If you use the hyperref package, please uncomment the following two lines
% to display URLs in blue roman font according to Springer's eBook style:
%\usepackage{color}
%\renewcommand\UrlFont{\color{blue}\rmfamily}
%\urlstyle{rm}
%

\usepackage[style=apa]{biblatex} % <- AGREGAR ESTA LÍNEA
\addbibresource{thebiblio.bib} % <- AGREGAR ESTA LÍNEA
\begin{document}
%
% ENGLISH VERSION
\title{Anomaly Detection in biomedical data by means of Autoencoders}

%\titlerunning{Abbreviated paper title}
% If the paper title is too long for the running head, you can set
% an abbreviated paper title here
%
\author{Matías Cisnero\inst{1} \and
Juliana Gambini\inst{1,2}\orcidID{0000-0002-4534-1402}}
%
\authorrunning{F. Author et al.}
% First names are abbreviated in the running head.
% If there are more than two authors, 'et al.' is used.
%
\institute{LIDEC, Universidad Nacional de Hurlingham, Pcia. de Buenos Aires, Argentina
\email{juliana.gambini@unahur.edu.ar} \and
CPSI, Universidad Tecnológica Nacional, Regional Buenos Aires}
%
\maketitle              % typeset the header of the contribution
%
\begin{abstract}
The abstract should briefly summarize the contents of the paper in
150--250 words. This is the English version of the abstract that 
provides a comprehensive overview of the research conducted, 
the methodology used, and the main findings obtained.

% keywords in lowercase
\keywords{Anomaly detection\and Biomedical Data\and Autoencoders}
\end{abstract}

% SPANISH VERSION
% Clear the previous title info and set Spanish version

% Spanish title and abstract
\begin{center}
\vspace{1cm}
{\Large \bfseries\boldmath Título de la contribución}

\vspace{0.5cm}
\end{center}

\begin{abstract}
El resumen debe resumir brevemente el contenido del artículo en
150-250 palabras. Esta es la versión en español del resumen que
proporciona una visión general completa de la investigación realizada,
la metodología utilizada y los principales hallazgos obtenidos.

% palabras clave en minúsculas
\keywords{primera palabra clave\and segunda palabra clave\and otra palabra clave}
\end{abstract}


\vspace{0.5cm}
\setcounter{page}{1}

%
\section{Introducción}

Motivación, hipótesis y estado del arte

\section{Materiales y Métodos}
\subsection{Conjunto de datos}
\subsection{Autoencoders}
\section{Resultados}
\section{Conclusiones y trabajos futuros}

%
% ---- Bibliography ----
%

\printbibliography

\end{document}